% Independently audited by root, four_regular, and path_stability.
\section{Adaptive attachments at the full pancyclicity threshold}
\label{sec:rhost-full-threshold}

The endpoint failure probability need not be small enough for a direct union bound.
We instead select the attachment edges after exposing each connector tree.
Long base cycles supply spare candidate edges.
For the short cycles, the number of candidates grows with their length.

\begin{theorem}[The leading term at the full threshold]\label{thm:rhost-full-threshold}
Suppose
\[
 np-\ln n-\ln\ln n\longrightarrow\infty.
\]
Then with high probability
\[
 \operatorname{Pex}(G(n,p))=(1+o(1))\log_2 n.
\]
\end{theorem}

\begin{proof}
Put
\[
 w=\min\{np-\ln n-\ln\ln n,\ln\ln n\},\qquad D=w/100000.
\]
Expose five independent graphs with
\[
 p_1=p_2=p_3=p_5=D/n,\qquad
 p_4=(\ln n+\ln\ln n+w/2)/n.
\]
Their union is dominated by $G(n,p)$.
We have $D\to\infty$, without any specified rate.

Use
\[
 \beta=20\ln D/D,\quad a=\lfloor(5\beta)^{-1}\rfloor,\quad
 M=\lfloor n/1000\rfloor,\quad d=\lceil\log_a M\rceil.
\]
Thus $d\sim\ln n/\ln D=o(\ln n)$.
Lemmas~\ref{lem:rhost-threshold-expansion} and~\ref{lem:rhost-threshold-trees} supply the robust expanders and equal-depth leaf trees.
In particular, every permitted reservoir contains an induced expander of size $n-o(n)$.
It contains the required rooted trees and two-tree structures.
Each leaf tree has exactly $M$ leaves at depth $d$ and fewer than $3M$ vertices.
Write $H=\lfloor n/4000\rfloor$ and $\delta=\exp(-D/5000)=o(1)$.
A fresh endpoint has a neighbour in any specified $H$-vertex subset with failure at most $\delta$.
The same bound holds for a specified leaf set after at most $2H$ earlier endpoint queries.

Set
\[
 K_0=\lceil\log_2(4d)\rceil,\quad
 \ell_*=(K_0+2)(d+3),\quad b=\lceil\sqrt{2\ell_*}\rceil,
\]
\[
 c=8b+5,\qquad T=c+4d+3=4d+8b+8,
\]
\[
 K=\left\lfloor\log_2\frac{nw}{100\ln n}\right\rfloor,
 \qquad L_0=2^{K+1}-1,\qquad m=\lfloor n2^{-K}\rfloor.
\]
These parameters are fixed before exposure.
They satisfy
\[
 T,b,d,m=o(\ln n),\quad \ell_*=o(L_0),\quad L_0\gg n^{0.9},
 \quad 2L_0\le\frac{nw}{25\ln n}.
\]
Also, $\ell_*/K\to\infty$, since $D\le\ln\ln n$ and $\ln d\sim\ln\ln n$.

\emph{Initial cycles.}
In $G_1$, select vertex-disjoint bit cycles $C_i$ of lengths $2^i+2$, for $0\le i\le K_0$.
Select a digit base cycle $C_{\rm short}$ of length $c$.
Select one further cycle $S_j$ of each length $3\le j\le T$.
Lemma~\ref{lem:rhost-cycle-list} applies: all lengths are $O(\ln n)$, and their multiplicities are at most three.
The selected cycles use $O(T^2+d)$ vertices.
Also select an edge $f$ outside their vertices.
Such an edge exists with high probability: $G_1$ has $\Omega(nD)$ edges and maximum degree at most $(\ln n)^2$,
while the selected vertices number only $O((\ln n)^2)$.
Let $H_1$ contain all these cycles and $f$.

\emph{An adaptive chain of cycles.}
We describe a connector procedure in $G_2$.
Expose first only edges with both endpoints outside $H_1$.
The universal robust-expansion event then determines every required reservoir expander.
Edges from vertices of $H_1$ to the reservoir remain independent until queried.

Start at a specified endpoint outside the reservoir.
Query at most $H$ fresh candidate neighbours in the current expander.
If a neighbour $z$ exists, embed a leaf tree rooted at $z$.
This first step fails with probability at most $\delta$.
Suppose the next object is an untouched cycle $Q$ of length $j$.
Take a matching of $\lfloor j/2\rfloor$ candidate edges in $Q$.
For a candidate edge $xy$, first query the edges from $x$ to the previous leaf set.
If one exists, it gives an incoming path of length $d+2$ from the preceding endpoint to $x$.
Temporarily exclude this path's internal vertices from the next reservoir.
Then query at most $H$ fresh edges from $y$ to the resulting expander.
If this succeeds, retain the incoming path, select $xy$, and embed the next leaf tree at the chosen neighbour of $y$.
If either query fails, discard that candidate edge and retain no part of its tentative incoming path.
The previous leaf tree remains available for the next candidate.

Each candidate has two endpoints unused in any previous candidate.
The reservoir was exposed before all these endpoint queries.
Even after conditioning on earlier failures and on the tentative incoming path, the edges from the new endpoint $y$ remain fresh.
Thus each conditional candidate failure probability is at most $2\delta$.
A cycle of length $j$ fails with probability at most $(2\delta)^{\lfloor j/2\rfloor}$.
For any list with bounded multiplicity of cycle lengths, their total failure probability is at most
\[
 O\!\left(\sum_{j\ge3}(2\delta)^{\lfloor j/2\rfloor}\right)=o(1).
\]
At the end, one fixed target endpoint attaches to the final leaf set with failure at most $\delta$.
Only the chosen paths, not the unused tree vertices or failed queries, enter the constructed subgraph.

Apply this procedure first to the bit cycles, starting at one endpoint of $f$ and ending at the other.
The selected edge of $C_i$ is denoted $e_i$.
Together with $f$ and the $K_0+2$ connectors, these edges form a cycle $C_*$ of length $\ell_*$.
Each other arc of $C_i$ replaces $e_i$ with gain exactly $2^i$.

Fix an edge $e_{\rm short}$ of $C_{\rm short}$.
Apply the procedure again to $S_3,\ldots,S_T$.
Start at the endpoint of $f$ used for the first bounded query, and end at one endpoint of $e_{\rm short}$.
The starting endpoint has only $H$ earlier queries, so another $H$ fresh candidate edges remain available.
This second chain is disjoint from the first construction except at its designated endpoint.
It contains all single-cycle blocks needed for the short lengths.

\emph{Adaptive binary and digit detours.}
Expose the $G_3$ reservoir edges outside the entire graph constructed so far.
All its base-cycle endpoint edges remain fresh.
For each $i>K_0$, embed two depth-$d$ leaf trees whose roots are joined by a stem of length
\[
 2^i+1-2d-2=2^i-2d-1>0.
\]
The total tree size is below $n/50$.
Choose an unused candidate edge $xy$ of $C_*$, and query its two endpoints against the two leaf sets.
If both attachments succeed, the resulting detour has length $2^i+1$.
Its replacement gain over $xy$ is $2^i$.
If an attachment fails, abandon that candidate edge and try another with two fresh endpoints.

Preselect $2(K-K_0)$ pairwise vertex-disjoint candidate edges of $C_*$,
avoiding $f$ and all endpoints of $e_0,\ldots,e_{K_0}$.
There are enough because $\ell_*/K\to\infty$.
Every conditional failure probability is at most $2\delta$.
The expected number of failures in this entire candidate pool is at most $4\delta(K-K_0)$.
If the pool is exhausted before completing all detours, at least $K-K_0$ candidates have failed.
Markov's inequality shows that this event has probability at most $4\delta=o(1)$.
After each success, retain only its detour path and re-form the next reservoir.

The digit base cycle has at least $4b$ pairwise vertex-disjoint candidate edges avoiding the endpoints of $e_{\rm short}$.
Use them to obtain $2b$ successful detours, in two groups of $b$.
The desired detour lengths are
\[
 2d+3+j,\qquad 2d+3+jb\qquad(0\le j<b).
\]
The two-tree stems have lengths $1+j$ and $1+jb$, so they are positive.
The same expected-failures argument succeeds with probability $1-o(1)$.
The designated edges need not be consecutive on the base cycle.
Choosing one detour from each group gives every cycle length in
\[
 [T+1,T+b^2].
\]
Indeed, the length is $c-2+(2d+3+j_0)+(2d+3+bj_1)=T+1+j_0+bj_1$.
This interval overlaps $[\ell_*,\ell_*+L_0]$, since $T=o(\ell_*)$ and $b^2\ge2\ell_*$.

\emph{Hamilton completion and final coverage.}
Delete $f$, $e_{\rm short}$, and the selected edge of each $S_j$ when choosing a spanning path through the constructed blocks.
Use every bit detour, every digit detour, and the long arc of each $S_j$ in this path.
The separate short-cycle chain joins one end of the binary path to one end of the digit path.
Thus all constructed vertices lie on one spanning path.
The binary path is a contiguous segment of it.
All deleted edges remain present as chords in the constructed subgraph.

The total used vertex count is at most $2L_0$.
The binary detours contribute $L_0+O(\ell_*)$ vertices.
The short-cycle chain and digit gadget contribute only a polynomial in $\ln n$.
For the remaining vertex set $X$,
\[
 |X|p_4-\ln|X|-\ln\ln|X|
 \ge(1/2-1/25-o(1))w\longrightarrow\infty.
\]
Lemma~\ref{lem:rhost-threshold-hamilton} therefore joins the two remaining endpoints through all vertices of $X$.
This is the completion input from Alon--Krivelevich Lemma~2.1 and Claim~4.7.
Its endpoint failure probability is $o(1)$.
This gives a Hamilton cycle.

The bit cycles give every length in $[\ell_*,\ell_*+L_0]$.
Replacing bit detours on the Hamilton cycle gives every length in $[n-L_0,n]$.
The individual cycles $S_j$ and the digit gadget cover all lengths below the first interval.
Use $G_5$ only for the long interval-covering chords.
Lemma~\ref{lem:rhost-threshold-windows} applies because the binary segment is contiguous, has length $L_0+\ell_*-1$,
and $\ell_*=o(n^{0.9})=o(L_0)$.
Each required window has at least $n^{1.8}/100$ candidate edges.
The failure probability is at most $m\exp(-Dn^{0.8}/100)=o(1)$.
Their overlapping length intervals cover the remaining gap.

The constructed graph has excess at most
\[
 K+2b+T+m+O(1)=\log_2 n+o(\ln n).
\]
The short-cycle chain contributes $T-2$ to cycle rank.
Failed queries contribute no retained edges.
The universal cycle-count lower bound gives the asserted equality to leading order.
\end{proof}

\begin{corollary}[A quantitative full-threshold error]\label{cor:rhost-full-threshold-error}
Write $\omega=np-\ln n-\ln\ln n\to\infty$ and $w=\min\{\omega,\ln\ln n\}$.
Then with high probability
\[
 \operatorname{Pex}(G(n,p))
 \le\log_2n+(4+o(1))\frac{\ln n}{\ln w}.
\]
\end{corollary}

\begin{proof}
The preceding construction has $\ln a=(1+o(1))\ln w$, so $d\sim\ln n/\ln w$.
Moreover, $b/d=O(\sqrt{K_0/d})=o(1)$ and $m/d=O(\ln w/w)=o(1)$.
Therefore $T=4d+o(d)$, and the displayed excess bound gives the result.
\end{proof}
